% ---- IEEE Title Block ----
\title{Image Enhancement with Deep Learning:\\A Comparative Study of CNN-Based Image Restoration Methods}

\author{
\IEEEauthorblockN{Giuseppe Bellamacina, Daniele Barbagallo, Salvo Iurato, Mattia Campanella}
\IEEEauthorblockA{
Department of Mathematics and Computer Science\\
University of Catania, Italy\\
Deep Learning Course --- A.Y. 2025/2026\\
Professors: G.M. Farinella, F. Ragusa
}
}

\maketitle

\begin{abstract}
We present a comprehensive comparative study of deep learning architectures for image restoration under severe Gaussian noise ($\sigma = 100$). Eight distinct architectures are implemented within a unified, modular framework: UNet, UNet Residual, Attention UNet, DnCNN, Denoising Autoencoder, RIDNet, NAFNet-V2, and Pix2Pix (conditional GAN). All models are trained on the DIV2K dataset using consistent degradation and evaluation protocols. Our experiments demonstrate that NAFNet-V2 with a multi-stage training pipeline achieves the best overall performance (PSNR = 26.36~dB, SSIM = 0.747), followed by Pix2Pix with residual learning (PSNR = 26.08~dB) and RIDNet (PSNR = 25.83~dB). We perform extensive ablation studies on architectural choices including upsampling strategies, attention mechanisms, loss functions, GAN stabilization, and multi-stage training. Results show that residual learning, advanced loss functions (Charbonnier), and staged training strategies are critical factors for high-quality denoising under heavy noise conditions.
\end{abstract}

\begin{IEEEkeywords}
Image denoising, convolutional neural networks, residual learning, attention mechanisms, image restoration, deep learning
\end{IEEEkeywords}
